% Generated draft: requires mathematical and source review.
\documentclass[11pt]{article}
\usepackage[margin=25mm]{geometry}
\usepackage{fontspec}
\setmainfont{Latin Modern Roman}
\usepackage{amsmath,amssymb,mathtools,unicode-math}
\setmathfont{Latin Modern Math}
\usepackage{enumitem,needspace,xurl,hyperref}
\setlength{\parindent}{0pt}
\setlength{\parskip}{5pt}
\setlength{\emergencystretch}{3em}
\newcommand{\N}{\mathbb{N}}
\newcommand{\Z}{\mathbb{Z}}
\newcommand{\Q}{\mathbb{Q}}
\newcommand{\R}{\mathbb{R}}
\newcommand{\C}{\mathbb{C}}
\newcommand{\F}{\mathbb{F}}
\newcommand{\A}{\mathbb{A}}
\newcommand{\PP}{\mathbb{P}}
\newcommand{\E}{\mathbb{E}}
\newcommand{\eps}{\varepsilon}
\newcommand{\dd}{\,\mathrm{d}}
\begin{document}
\section*{UnsolvedMath ID: 30002178}
\textbf{Source problem code:} catalogue entry
\subsection{Definitions and mathematical statement}
Let \(\N_{>0}=\{1,2,3,\ldots\}\). For \(d\in\N_{>0}\), define the set of \(d\)-th roots of unity by
\[
C_d=\{e^{2\pi i k/d}:k\in\Z\}\subset\C.
\]
For \(m\in\N_{>0}\), its \(m\)-fold Minkowski sum is
\[
mC_d=\{z_1+\cdots+z_m:z_1,\ldots,z_m\in C_d\},
\]
where repetitions are allowed. For a finite set \(A\subset\C\) containing a nonzero element, put
\[
\min_0(A)=\min\{|z|:z\in A\setminus\{0\}\}.
\]
This minimum is defined for \(mC_d\), since this set is finite and contains the nonzero number \(m\).

The original source record asks for the following asymptotic lower bound:
\[
\forall m\in\N_{>0}\ \exists E(m)\in\R_{>0}\ \exists D(m)\in\N_{>0}\ \forall d\in\N_{>0},\qquad
 d\ge D(m)\ \Longrightarrow\ \min_0(mC_d)>d^{-E(m)}.
\]
Equivalently, for each fixed \(m\), every nonzero sum of \(m\) roots in \(C_d\) should have modulus greater than \(d^{-E(m)}\) for all sufficiently large positive integers \(d\), with the exponent independent of the chosen roots and of \(d\).

\textbf{Unresolved boundary qualification.} The supplied clean statement instead asks for
\[
\forall m\in\N_{>0}\ \exists E(m)\in\R_{>0}\ \forall d\in\N_{>0},\qquad
\min_0(mC_d)>d^{-E(m)}.
\]
This differs from the original wording ``as \(d\to\infty\).'' Moreover, at \(m=d=1\), the latter displayed inequality reads \(1>1\). Thus the all-positive-\(d\) formulation cannot be retained as an equivalent formulation without a boundary correction. The supplied evidence does not determine whether the intended correction is an eventual bound or an explicit restriction on \(d\); both supplied formulations and this discrepancy require human review.
\subsection{Short English statement}
For each fixed number \(m\) of summands, the original question asks whether every nonzero sum of \(m\) \(d\)-th roots of unity has modulus greater than a fixed inverse power of \(d\) for all sufficiently large \(d\). The supplied clean version requires the inequality for every positive \(d\), creating an unresolved boundary discrepancy at \(m=d=1\).
\subsection{Sources}
\begin{itemize}
\item Local registry: \texttt{data/unsolvedmath/problems.json}, supplied record with permanent UnsolvedMath ID \(30002178\), title \emph{Lower Bounds for Sums of Roots of Unity}, exact source problem code \texttt{catalogue entry}, in the Oberwolfach Reports Open Problems collection. No individual proposer is supplied; the proposed year is recorded as 2012.
\item Supplied source citation: \emph{Low-Dimensional Topology and Number Theory} (2013), Oberwolfach Reports, pp.\ 2541--2596, DOI \texttt{10.4171/owr/2012/42}. The supplied literature list identifies Oberwolfach Report 42, Conjecture 7: \url{https://doi.org/10.4171/owr/2012/42}.
\item Supplied literature reference: Barber, small sums of five roots of unity, \url{https://doi.org/10.1112/blms.12826}.
\item Supplied literature reference: Kalogirou, small sums of roots of unity, \url{https://arxiv.org/abs/2607.06098}.
\item Literature assessment dated exactly \texttt{2026-08-21}: the supplied curation labels the general fixed-\(m\) lower-bound problem open and reports that recent work provides stronger constructions of small sums and isolated fixed-conductor results, without establishing the conjectured uniform lower bound. This is supplied dated curation, not an independently verified literature assessment.
\item The registry labels the statement \texttt{corrected\_verified} and says that it was checked against the official report. These are supplied verification claims; the report and cited literature have not been independently consulted for this draft.
\end{itemize}
\end{document}
